
\subsection{Performa Deteksi dan Segmentasi (YOLOv8)}
\label{sec:yolo_performance}

Analisis performa deteksi dan segmentasi meninjau dua aspek: konfigurasi pelatihan teknis dan hasil evaluasi metrik kuantitatif. Evaluasi ini bertujuan memvalidasi kemampuan model dalam mengisolasi area lubang jalan (\textit{Region of Interest}) secara akurat sebelum diproses lebih lanjut.

\subsubsection{Konfigurasi Pelatihan Model}

Berdasarkan eksperimen awal, arsitektur \textbf{YOLOv8m-seg} (\textit{Medium}) dipilih menggantikan varian Nano. Varian Medium menawarkan keseimbangan yang lebih baik antara kedalaman ekstraksi fitur dan latensi inferensi, yang krusial untuk menangani tekstur jalan yang kompleks.

Pelatihan dilakukan menggunakan akselerator \textbf{NVIDIA GeForce RTX 4060 Ti} (16GB VRAM) dengan kerangka kerja PyTorch 2.5.1. Parameter pelatihan utama meliputi:
\begin{itemize}
    \item \textbf{Resolusi Input:} $960 \times 960$ piksel. Resolusi tinggi ini dipilih untuk meningkatkan daya tangkap terhadap lubang berukuran kecil atau yang berada jauh dari kamera.
    \item \textbf{Durasi \& Strategi:} Total 192 \textit{epoch} dengan mekanisme \textit{early stopping} (patience=100). Konvergensi optimal tercapai dan model terbaik disimpan pada \textit{epoch} ke-92.
    \item \textbf{Hiperparameter:} Ukuran \textit{batch} 8, pengoptimal SGD (\textit{Stochastic Gradient Descent}) dengan \textit{learning rate} awal 0,005, momentum 0,937, dan \textit{weight decay} 0,0005.
\end{itemize}

Selama pelatihan, kurva fungsi kerugian (\textit{loss}) untuk klasifikasi, regresi \textit{box}, dan segmentasi menunjukkan tren penurunan yang stabil (Gambar~\ref{fig:training_loss}), mengindikasikan model berhasil mempelajari fitur representatif tanpa mengalami \textit{overfitting} yang signifikan.

\begin{figure}[htbp]
  \centering
  \framebox{\parbox{0.8\textwidth}{\centering
    \vspace{1cm}
    \textbf{[Placeholder Grafik Loss]} \\
    \small\textit{Masukkan grafik training/validation loss dari file results.png di sini.}
    \vspace{1cm}
  }}
  \caption{Kurva penurunan loss pelatihan dan validasi model YOLOv8m-seg.}
  \label{fig:training_loss}
\end{figure}

\subsubsection{Evaluasi Metrik Deteksi}

Evaluasi kuantitatif dilakukan pada dataset validasi berisi 300 citra dengan total 565 instansi objek. Tabel~\ref{tab:yolo_metrics} merangkum kinerja model. Secara agregat, model mencapai \textit{Mean Average Precision} (mAP) pada ambang batas IoU 0,5 sebesar \textbf{0,789}.

\begin{table}[h]
\caption{\label{tab:yolo_metrics}Metrik Evaluasi Segmentasi YOLOv8m-seg pada Dataset Validasi.}
\begin{center}
\begin{tabular}{lcccc}
\toprule
Kelas & Precision (P) & Recall (R) & mAP@0.5 & mAP@0.5:0.95 \\
\midrule
\textbf{Semua Kelas} & \textbf{0.765} & \textbf{0.756} & \textbf{0.789} & \textbf{0.619} \\
Pothole & 0.795 & 0.912 & 0.892 & 0.685 \\
Manhole & 0.735 & 0.600 & 0.687 & 0.552 \\
\bottomrule
\end{tabular}
\end{center}
\end{table}

Analisis mendalam terhadap hasil tersebut menunjukkan:
\begin{enumerate}
    \item \textbf{Sensitivitas Tinggi:} Skor \textit{Recall} sebesar \textbf{0,912} pada kelas \textit{Pothole} mengindikasikan bahwa sistem mampu mendeteksi lebih dari 91\% lubang jalan yang ada. Hal ini meminimalkan risiko \textit{false negatives} yang berbahaya bagi keselamatan pengguna jalan.
    \item \textbf{Presisi Segmentasi:} Nilai mAP@0.5 sebesar \textbf{0,892} menegaskan bahwa mask segmentasi sangat presisi mengikuti kontur kerusakan yang ireguler, yang vital untuk akurasi pengukuran diameter di tahap selanjutnya.
    \item \textbf{Diskriminasi Kelas:} Meskipun skor \textit{Manhole} lebih rendah (mAP 0,687) akibat ketidakseimbangan data, model tetap efektif membedakan lubang berbahaya dari infrastruktur jalan normal.
\end{enumerate}

Gambar~\ref{fig:detection_results} memperlihatkan contoh kualitatif hasil segmentasi pada kondisi pencahayaan berbeda, menunjukkan kemampuan model beradaptasi terhadap variasi visual lingkungan.

\begin{figure}[htbp]
  \centering
  \framebox{\parbox{0.8\textwidth}{\centering
    \vspace{1.5cm}
    \textbf{[Placeholder Hasil Deteksi]} \\
    \small\textit{Masukkan contoh gambar hasil deteksi (val\_batch0\_pred.jpg) di sini.}
    \vspace{1.5cm}
  }}
  \caption{Visualisasi hasil deteksi dan segmentasi pada data uji.}
  \label{fig:detection_results}
\end{figure}


